\chapter{相对论量子力学}

\section{Dirac方程}
曾谨言——《量子力学 II》第五版，Eq. 11.2.24，证明
\begin{proof}
    \begin{equation*}
        \left[ \frac{1}{2} \bSigma, H \right] =
                        \left[ \frac{1}{2}(\Sigma_x + \Sigma_y + \Sigma_z), 
                        c(\alpha_x p_x + \alpha_y p_y + \alpha_z p_z) \right]
    \end{equation*}
    只选$\Sigma_x$这一项进行计算，有以下
    \begin{equation}
        \begin{aligned}
            \left[\frac{1}{2}\Sigma_x, c(\alpha_x p_x + \alpha_y p_y + \alpha_z p_z)\right]
          &= \left[ \frac{1}{2} \Sigma_x, c(\alpha_y p_y + \alpha_z p_z) \right] \\
          &=    c \left[ \frac{1}{2} \Sigma_x, \alpha_y p_y\right] 
              + c \left[ \frac{1}{2} \Sigma_x, \alpha_z p_z\right]   \\
          &=    c \left[ \frac{1}{2} \Sigma_x, \alpha_y\right] p_y
              + c \left[ \frac{1}{2} \Sigma_x, \alpha_z\right] p_z \\
          &=    {\rm i} c \alpha_z p_y  - {\rm i} c \alpha_y p_z \\
          &= -ic(\alpha_y p_z - \alpha_z p_y)
        \end{aligned}
    \end{equation}
    按上述分别对$\Sigma_y$、$\Sigma_z$进行处理，最后可得
    \begin{equation*}
        \left[ \frac{1}{2} \bSigma, H \right] = -{\rm i} c (\bm{\alpha} \times \bm{p})
    \end{equation*}
\end{proof}

曾谨言——《量子力学 II》第五版，练习1，Eq. 11.2.34
\begin{proof}
    \begin{equation}
        \begin{aligned}
            (\bSigma \cdot \bmA) (\bSigma \cdot \bmB) =&
            \begin{pmatrix}
                \bsigma \cdot \bmA  &  0 \\
                0                         &  \bsigma \cdot \bmA
            \end{pmatrix}
            \begin{pmatrix}
                \bsigma \cdot \bmB  &  0 \\
                0                         &  \bsigma \cdot \bmB
            \end{pmatrix} \\
        =& \begin{pmatrix}
            (\bsigma \cdot \bmA) (\bsigma \cdot (\bmB)  &  0 \\
            0  &  (\bsigma \cdot \bmA) (\bsigma \cdot (\bmB) \\
           \end{pmatrix} \\
        =& \begin{pmatrix}
            \bmA \cdot \bmB + {\rm i} \bsigma \cdot (\bm{A} \times \bmB)  &  0 \\
            0  &  \bmA \cdot \bmB + {\rm i} \bsigma \cdot (\bm{A} \times \bmB) \\
           \end{pmatrix} \\
        =& \bmA \cdot \bmB + {\rm i} \bSigma (\bm{A} \times \bmB)
        \end{aligned}
    \end{equation}
    最后一个等号的第一项其实是$\bm{I} \cdot (\bmA \cdot \bmB)$，$\bm{I}$是一个与
    $\bSigma$相同维度的单位矩阵。
\end{proof}

曾谨言——《量子力学 II》第五版，练习1，Eq. 11.2.35
\begin{proof}
    \begin{equation}
        \begin{aligned}
            (\balpha \cdot \bmA) (\balpha \cdot \bmB) =&
            \begin{pmatrix}
                0                   &  \bsigma \cdot \bmA  \\
                \bsigma \cdot \bmA   & 0
            \end{pmatrix}
            \begin{pmatrix}
                0                   &  \bsigma \cdot \bmB  \\
                \bsigma \cdot \bmB   & 0
            \end{pmatrix} \\
        =& \begin{pmatrix}
            (\bsigma \cdot \bmA) (\bsigma \cdot \bmB)  &  0 \\
            0  &  (\bsigma \cdot \bmA) (\bsigma \cdot \bmB) \\
           \end{pmatrix} \\
        =& \bmA \cdot \bmB + {\rm i} \bSigma (\bm{A} \times \bmB)
        \end{aligned}
    \end{equation}
    最后一个等号的第一项其实是$\bm{I} \cdot (\bmA \cdot \bmB)$，$\bm{I}$是一个与
    $\bSigma$相同维度的单位矩阵。
\end{proof}

曾谨言——《量子力学 II》第五版，练习1，Eq. 11.2.36
\begin{proof}
    首先，我们有
    \begin{equation}
            (\balpha \cdot \bmA) (\bSigma \cdot \bmB) =
            \begin{pmatrix}
                0                   &  \bsigma \cdot \bmA  \\
                \bsigma \cdot \bmA   & 0
            \end{pmatrix}
            \begin{pmatrix}
                \bsigma \cdot \bmB   & 0  \\
                0  &  \bsigma \cdot \bmB 
            \end{pmatrix}
        = \begin{pmatrix}
            0  &  (\bsigma \cdot \bmA) (\bsigma \cdot \bmB) \\
            (\bsigma \cdot \bmA) (\bsigma \cdot \bmB) &  0  \\
           \end{pmatrix} \\
    \end{equation}
    同样的，有
    \begin{equation}
            (\bSigma \cdot \bmA) (\balpha \cdot \bmB) =
            \begin{pmatrix}
                \bsigma \cdot \bmA   & 0  \\
                0  &  \bsigma \cdot \bmA 
            \end{pmatrix}
            \begin{pmatrix}
                0                   &  \bsigma \cdot \bmB  \\
                \bsigma \cdot \bmB   & 0
            \end{pmatrix}
        = \begin{pmatrix}
            0  &  (\bsigma \cdot \bmA) (\bsigma \cdot \bmB) \\
            (\bsigma \cdot \bmA) (\bsigma \cdot \bmB) &  0  \\
           \end{pmatrix} \\
    \end{equation}
    由以上两式可知$(\balpha\cdot\bmA)(\bSigma\cdot\bmB)=(\bSigma\cdot\bmA)(\balpha\cdot\bmB)$。

    现在，我们对上面最后一个等式做进一步计算，有
    \begin{equation}
        \begin{aligned}
            \begin{pmatrix}
                0  &  (\bsigma \cdot \bmA) (\bsigma \cdot \bmB) \\
                (\bsigma \cdot \bmA) (\bsigma \cdot \bmB) &  0  \\
            \end{pmatrix}
            =& \begin{pmatrix}
                0 & \bmA \cdot \bmB + {\rm i} \bsigma \cdot (\bm{A} \times \bmB) \\
                \bmA \cdot \bmB + {\rm i} \bsigma \cdot (\bm{A} \times \bmB) & 0 \\
            \end{pmatrix} \\
            =& \begin{pmatrix}
                0 & \bmA \cdot \bmB \\
                \bmA \cdot \bmB & 0 \\
            \end{pmatrix} 
            + \begin{pmatrix}
                0 & {\rm i} \bsigma \cdot (\bm{A} \times \bmB) \\
                {\rm i} \bsigma \cdot (\bm{A} \times \bmB) & 0 \\
            \end{pmatrix} \\
            =& - r_5 \bmA \cdot \bmB + i \balpha \cdot (\bmA \times \bmB)
        \end{aligned}
    \end{equation}
    其中
    \begin{equation*}
        r_5 = \begin{pmatrix}
            0 & -I \\
            -I & 0
        \end{pmatrix}
        = -\begin{pmatrix}
            0 & I \\
            I & 0
        \end{pmatrix}
    \end{equation*}
\end{proof}

曾谨言——《量子力学 II》第五版，练习2
\begin{proof}
    此处$\bsigma \cdot \bmp / |p|$应为$\bSigma\cdot\bmp/|p|$，是一个维度为4的矩阵，要证明
    它是守恒量，需要证明其与哈密顿量对易：
    \begin{equation}
        \begin{aligned}
            \left[ \frac{\bSigma\cdot\bmp}{|p|}, \hat{\bmH} \right] =&
            \left[ \frac{\bSigma\cdot\bmp}{|p|}, c\balpha\cdot\bmp \right]
            = \frac{1}{|p|} \left[ \sum_i\Sigma_i p_i, \sum_j\alpha_j p_j \right] \\
            =& \frac{1}{|p|} \sum_{ij} \left[ \Sigma_i p_i, \alpha_j p_j \right] \\
            =& \frac{1}{|p|} \sum_{ij} 
               \left( \Sigma_i [p_i, \alpha_j p_j] + [\Sigma_i, \alpha_j p_j] p_i \right) \\
            =& \frac{1}{|p|} \sum_{ij} \left( 
                \cancel{\Sigma_i \alpha_i[p_i, p_j]} + \cancel{\Sigma_i [p_i, \alpha_j]p_j}
              + \cancel{\alpha_i[\Sigma_i, p_j]p_i} + [\alpha_i, \Sigma_i]p_j p_i
               \right) \\
            =& \frac{2{\rm i}}{|p|}\sum_{ijk}\epsilon_{ijk}\alpha_k p_j p_i \\
            =& \frac{2{\rm i}}{|p|}\bmp \times \bmp \cdot \balpha = 0
        \end{aligned}
    \end{equation}
    第四个等号右边第二三项能消去是由于动量算符与自旋算符的对易关系。\uline{其本征值待求}。
\end{proof}
